% Exported from Math++. Compile with pdflatex.
\documentclass[11pt,a4paper]{article}
\usepackage[T1]{fontenc}
\usepackage[utf8]{inputenc}
\usepackage{lmodern}
\usepackage{amsmath,amssymb,amsthm,mathtools}
\usepackage{graphicx,booktabs,xcolor}
\usepackage[hidelinks]{hyperref}
\usepackage[margin=25mm]{geometry}
\theoremstyle{plain}
\newtheorem{theorem}{Theorem}
\newtheorem{lemma}[theorem]{Lemma}
\newtheorem{corollary}[theorem]{Corollary}
\newtheorem{proposition}[theorem]{Proposition}
\theoremstyle{definition}
\newtheorem{definition}[theorem]{Definition}
\newtheorem{example}[theorem]{Example}
\theoremstyle{remark}
\newtheorem*{remark}{Remark}
\setlength{\parindent}{0pt}\setlength{\parskip}{0.6em}
\begin{document}
\section*{Matrices}

Write a matrix with \texttt{pmatrix}, then work with it.

\begin{align*}
  A \coloneqq  \begin{pmatrix} 1 & 2 \\ 3 & 4 \end{pmatrix}
\end{align*}
\begin{align*}
  B \coloneqq  \begin{pmatrix} 0 & 1 \\ 1 & 0 \end{pmatrix}
\end{align*}
The product keeps the order you write it in:

\begin{align*}
  A B = \begin{pmatrix} 2 & 1 \\ 4 & 3 \end{pmatrix}
\end{align*}
\begin{align*}
  B A = \begin{pmatrix} 3 & 4 \\ 1 & 2 \end{pmatrix}
\end{align*}
\begin{align*}
  A^{-1} = \begin{pmatrix} -2 & 1 \\ \frac{3}{2} & -\frac{1}{2} \end{pmatrix}
\end{align*}
\begin{align*}
  A^T = \begin{pmatrix} 1 & 3 \\ 2 & 4 \end{pmatrix}
\end{align*}
\begin{equation*}
  \det A = -2
\end{equation*}
\begin{equation*}
  \operatorname{rank} A = 2
\end{equation*}
A matrix times a vector:

\begin{equation*}
  A \begin{pmatrix} 1 \\ 1 \end{pmatrix} = \begin{pmatrix} 3 \\ 7 \end{pmatrix}
\end{equation*}
\end{document}
